\chapter{Дүгнэлт}
\label{sec:conclusion}

Энэхүү судалгааны ажлаар Flutter дээрх харагдах хэсэг (frontend), Node.js/Express дээрх сервер талын хэсэг (backend), MongoDB өгөгдлийн сан болон OpenAI үйлчилгээг нэгтгэсэн ``AI Notebook'' системийг бодит түвшинд хэрэгжүүлж, баталгаажуулалтын урсгал, тэмдэглэлийн CRUD үйлдэл, хиймэл оюун ухаанаар стикер үүсгэх үндсэн боломжуудыг интеграцийн түвшинд амжилттай туршиж нотоллоо. Судалгааны үр дүнгээс үзэхэд модульчлагдсан архитектур, JWT-д суурилсан хамгаалалт, өгөгдлийн урсгал болон дарааллын диаграммын уялдаа нь системийн ажиллагааг ойлгомжтой, өргөтгөх боломжтой, инженерчлэлийн хувьд хэрэгжихүйц шийдэл болгож байгаа нь тогтоогдсон бөгөөд академик судалгааг бодит програм хангамжийн шийдэлтэй уялдуулсан практик ач холбогдолтой ажил болсныг харуулж байна.

\textbf{Цаашдын төлөвлөгөө:} Системийг дараагийн шатанд хөгжүүлэхдээ түлхэх мэдэгдлийн модулийг үйлдвэрлэлийн орчинд бүрэн нэвтрүүлж, токен шинэчлэх механизм, үүрэгт суурилсан эрхийн удирдлага, сүлжээ тасарсан үед ажиллах синхрончлолыг (offline-first synchronization) нэмэгдүүлэх, мөн AI функцийг тэмдэглэл хураангуйлах болон сэтгэл хөдлөлийн шинжилгээ зэрэг чиглэлээр өргөтгөх шаардлагатай. Үүнтэй зэрэгцэн Docker болон CI/CD-д суурилсан байршуулалт, мониторинг, аюулгүй байдлын гүнзгий баталгаажуулалт, хэрэглэгчийн бодит туршилтад тулгуурласан UX үнэлгээг үе шаттайгаар хэрэгжүүлснээр AI Notebook системийг үйлдвэрлэлийн хэрэглээнд бэлэн, тогтвортой платформ болгон хөгжүүлэх боломжтой.
